\chapter{Advanced Systems --- The Engineer's Model of the Machine}

Every chapter that follows rests on one shift in perspective: the C++ abstract machine is a fiction the compiler maintains for your convenience, and the real machine underneath --- pipelined, cached, virtually addressed, multi-core, talking to an OS across an expensive boundary --- behaves nothing like the simple sequential model in your head. This opening chapter establishes the mental model the rest of the volume sharpens.

\section*{Chapter Roadmap}
\begin{itemize}
\item 85.1 The Abstract Machine vs the Real Machine
\item 85.2 The Layers Between Source and Silicon
\item 85.3 Latency Numbers Every Engineer Must Know
\item 85.4 The Three Cost Models: Compute, Memory, and the OS Boundary
\item 85.5 The Volume's Throughline: Determinism and the Tail
\item 85.6 How to Read the Rest of This Volume
\end{itemize}

\section{The Abstract Machine vs the Real Machine}
C++ is defined against an abstract machine: sequential, in-order, every read returns the last write, operations cost what they look like. The compiler preserves only this model's observable behaviour (the as-if rule). The real machine violates the implied cost and ordering everywhere: out-of-order and speculative execution, caches (a miss $\sim$100$\times$ a hit), virtual addressing (TLB, page tables), many cores seeing each other's writes with delay, and an OS boundary costing hundreds of cycles.

\textbf{Why this matters.} The abstract machine is why code is correct; the real machine is why it is fast or slow. Hold both: reason about correctness in the abstract machine, performance in the real one. Almost every performance surprise is the gap between them.

\section{The Layers Between Source and Silicon}
\begin{lstlisting}[language=text, caption={The transformation stack; each layer can make identical source perform differently.}]
C++ source -> compiler front-end [Ch 81] -> optimizer/IR [Ch 80,89,104]
 -> machine code [Ch 89] -> CPU front-end (decode, predict) [Ch 86]
 -> out-of-order core [Ch 86] -> cache hierarchy + coherence [Ch 87,76]
 -> virtual memory (TLB, pages, NUMA) [Ch 88] -> OS boundary [Ch 98-101]
\end{lstlisting}

\textbf{Why this matters.} Performance problems live at a specific layer; fixing them requires identifying which. A mispredict is CPU front-end; a cache miss is memory hierarchy; a page fault is virtual memory; a \texttt{seq\_cst} stall is coherence; a hot-path syscall is the OS boundary. Mistaking the layer wastes effort.

\section{Latency Numbers Every Engineer Must Know}
\begin{itemize}
\item 1 cycle $\sim$0.3\,ns; L1 hit $\sim$1\,ns; branch mispredict $\sim$5\,ns; L2 $\sim$4\,ns; L3 $\sim$15\,ns.
\item DRAM $\sim$100\,ns ($\sim$300$\times$ a cycle); NVMe read $\sim$10--100\,$\mu$s; datacenter round trip $\sim$50--500\,$\mu$s.
\item HDD seek $\sim$10\,ms; cross-continent round trip $\sim$100+\,ms.
\end{itemize}

\textbf{Why this matters.} As ratios, these are the engineer's intuition: a cache miss costs $\sim$300 instructions (layout dominates compute); a syscall/round trip is so costly that avoiding it beats optimizing around it; a mispredict is cheap alone but lethal in a tight loop. Memorise the orders of magnitude; measure the exact nanoseconds (Chapters 101, 103).

\section{The Three Cost Models: Compute, Memory, and the OS Boundary}
\begin{enumerate}
\item \textbf{Compute} --- instruction count, pipelining, mispredicts, vectorization (Chapters 86, 89, 91--92).
\item \textbf{Memory} --- cache vs DRAM, prefetchability, layout, false sharing (Chapters 76, 79, 87--88, 90, 97).
\item \textbf{OS boundary} --- syscalls, I/O, scheduling, clocks (Chapters 96, 98--101).
\end{enumerate}

\textbf{Why this matters.} Identifying which model dominates is the first move in optimization. Numeric kernels are compute-bound (vectorize); graph traversals memory-bound (improve locality); network servers OS-boundary-bound (reduce syscalls). Applying the wrong model's techniques is the most common waste of effort (Chapters 103, 105 tell you which model you are in).

\section{The Volume's Throughline: Determinism and the Tail}
The target systems care less about average latency than about the tail: the 99.9th-percentile request, the worst case that violates an SLA.

\textbf{Why this matters.} Every villain here is a source of variance: page faults, GC pauses, lock convoys, cache misses, syscalls, mispredicts, thread migrations, slow-path allocations. The hot-path discipline (Chapter 106) eliminates jitter --- preallocate, pin, warm, avoid the kernel --- so the worst case approaches the average. Measure full distributions and tail percentiles, never means (Chapter 103).

\section{How to Read the Rest of This Volume}
Hardware (86--89); data \& vectorisation (90--92); concurrency (76--78, 93--96); memory management (79, 97); OS interface (98--102); discipline (80, 103--106).

\textbf{The discipline.} Read each chapter asking: What does the hardware or OS actually do? What does it cost (in the units of 85.3)? When should I not do this? The techniques are only as good as the judgement of when to apply them.
